\cleardoublepage
\chapter{METODOLOGI PENELITIAN}
\label{chap:metodologi-penelitian}
Jelaskan secara rinci metodologi penelitian yang digunakan dalam pelaksanaan tesis, termasuk pendekatan penelitian, metode pengumpulan data, teknik analisis data, dan langkah-langkah pelaksanaan penelitian. Metodologi penelitian ini harus disesuaikan dengan topik tesis yang dibahas. Berikut adalah contoh penjelasan metodologi penelitian yang dapat digunakan sebagai referensi (judul subbab dapat diubah sesuai dengan hal yang hendak dibahas dalam subbab tersebut).
Sebagai contoh, jika topik tesis adalah tentang masalah klasifikasi biji kopi berdasarkan citra kopi, maka metodologi penelitian yang digunakan dapat dijelaskan sebagai berikut:
\section{Pendekatan Penelitian}
Jelaskan pendekatan penelitian yang digunakan, misalnya pendekatan kualitatif, kuantitatif, atau campuran (\textit{mixed methods}). Jelaskan alasan pemilihan pendekatan tersebut dan bagaimana pendekatan tersebut sesuai dengan tujuan penelitian.
\section{Metode Pengumpulan Data}
Jelaskan metode pengumpulan data yang digunakan, misalnya wawancara, observasi, kuesioner, studi dokumentasi, atau eksperimen. Jelaskan bagaimana metode tersebut dilakukan, siapa yang menjadi responden atau partisipan, dan bagaimana data dikumpulkan secara sistematis.
\section{Teknik Analisis Data}
Jelaskan teknik analisis data yang digunakan, misalnya analisis statistik, analisis tematik, analisis isi, atau teknik analisis lainnya yang sesuai dengan jenis data yang dikumpulkan. Jelaskan bagaimana teknik analisis tersebut dilakukan dan bagaimana hasil analisis akan digunakan untuk menjawab pertanyaan penelitian.
\section{Langkah-langkah Pelaksanaan Penelitian}
Jelaskan langkah-langkah pelaksanaan penelitian secara rinci, mulai dari perencanaan penelitian, pengumpulan data, analisis data, hingga penarikan kesimpulan. Jelaskan juga jadwal pelaksanaan penelitian, sumber daya yang dibutuhkan, dan risiko yang mungkin dihadapi selama pelaksanaan penelitian beserta mitigasinya.    
\section{Jadwal Pelaksanaan Penelitian}
Jelaskan jadwal pelaksanaan penelitian secara rinci, mulai dari tahap perencanaan, pengumpulan data, analisis data, hingga penarikan kesimpulan. Jelaskan juga sumber daya yang dibutuhkan untuk setiap tahap penelitian dan risiko yang mungkin dihadapi selama pelaksanaan penelitian beserta mitigasinya. Jadwal pelaksanaan penelitian ini dapat disajikan dalam bentuk tabel atau diagram Gantt untuk memberikan gambaran yang jelas tentang alur pelaksanaan penelitian.  
